\documentclass[a4paper]{article}
\usepackage[T2A]{fontenc}
\usepackage[utf8]{inputenc}
\usepackage[russian]{babel}
\usepackage{geometry}
\usepackage{setspace} % Для управления интервалами
\usepackage{myrussian}

% Настройка полей согласно ГОСТ (приблизительно для диплома МГУ)
\geometry{
  left=3cm,
  right=1cm,
  top=2cm,
  bottom=2cm
}

% Убираем нумерацию страниц на титульном листе
\pagestyle{empty}

\begin{document}
\textbf{1.} Пусть в диаграмме узла \( K \) содержится \( n \) дуг.
Тогда каждая трехцветная раскраска представляет собой вектор из \( n
\) компонент, каждая в \( \Z^3 \). При этом, на каждом перекрестке на
раскраску накладывается ограничение \( 2b \equiv a + c\mod 3 \).
Получается, что каждая раскраска --- решение системы \( Mv = 0 \),
где \( M \) --- матрица \( n \times n \). Пусть размерность
пространства решений равна \( d \). Тогда количество решений, а
следовательно и количество раскрасок равно \( 3^d \).\\
Пусть число заузленности узла равно \( U(K) \). После изменения
одного из перекрестков изменяется одно из уравнений системы на
раскраски. Значит размерность пространства решений изменяется на \(
-1,0,1 \). После замены \( U(K) \) перекрестков по определению числа
заузленности получим тривиальный узел, у которого \( 3 \) раскраски и
размерность пространства решений равна \( 1 \).
Получаем, что у узла \( K \) размерность пространства решений \( d \) не
превышает \( U(K) + 1 \). Отсюда
\[
  c_3(K) \le 3^{U(K) + 1}.
\]
\textbf{3.} Будем рассматривать полином Джонса вместо скобки
Кауфмана, т.к. из формулы
\[
  X(L) = (-a)^{-3w(L)}\langle | L |\rangle
\]
следует, что скобка Кауфмана равна нулю \( \iff \) полином Джонса
равен нулю. Пусть \( L_+, L_-, L_0 \) - диаграммы узла \( L \) в
которых один перекресток заменен на положительный, отрицательный и
разведенный соответственно. Тогда полином Джонса удовлетворяет соотношению
\[
  a^4 X(L_+) - a^{-4}X(L_-) = (a^{-2} - a^2)X(L_0).
\]
Подставим \( a=1 \): \( X(L_+)(1) =  X(L_-)(1) \). То есть значение
полинома Джонса в единице не меняется при замене перекрестков. Так
как любой узел можно привести к тривиальному заменой перекрестков,
получем, что для произвольного зацепления значение полинома Джонса в
единице равно значению в единице полинома Джонса для тривиального зацепления.
Полином Джонса для тривиального зацепления из \( n \) компонент равен
\((-a^2 - a^{-2})^{n-1}  \). Получаем, что для любого зацепления из
\( n \) комепонент полином Джонса в единице равен \( -2^{n-1} \),
значит не равен тождественно нулю.

\textbf{4.} Группа кос из двух нитей состоит из одной образующей и
изоморфна группе \( \Z \). Группа кос из трех нитей --- группа с
двумя образующими \( \sigma_1, \sigma_2 \) с соотношением
\begin{equation}
  \label{braid3}
  \sigma_1
  \sigma _2 \sigma _1 = \sigma _2 \sigma _1 \sigma _2.
\end{equation}
Эта группа некоммутативна, т.к. из соотношения \ref{braid3} и
коммутативности \( \sigma _1 \sigma _2  = \sigma _2 \sigma _1\)
следует \( \sigma _1 = \sigma _2 \), что очевидно неверно. Получаем
неизоморфность групп.
\end{document}
